\section{Implementation and Mechanism Isolation}
\label{sec:implementation}

\subsection{Executable Selection Procedure}
\label{sec:impl:assessment}

Algorithm~\ref{alg:abstract} specifies the procedure that produces the
reported diagnosis-plus-gate results. The diagnostic report and validator
are deterministic. The language model supplies field values; the gate only
accepts or rejects proposed triples and cannot create missing facts.

\begin{algorithm}[t]
\caption{Source-grounded document-graph repair}
\label{alg:abstract}
\begin{algorithmic}[1]
\Require Input multiset $G$, source $X$, document node $h_0$, relation vocabulary $\mathcal R_X$
\Ensure Repaired graph $\widehat G$ and rejection log $J$
\State $G_s\gets P(G;h_0,\mathcal R_X)$
\State $D\gets D(G_s,X)$
\State $C\gets\textsc{Parse}(L_\theta(X,h_0,\mathcal R_X,G_s,D))$
\State $\widehat G\gets\emptyset$; $J\gets\emptyset$; $B\gets\emptyset$
\For{each distinct candidate $t=(h,r,o)$ in generation order}
    \If{$H(t)R(t)S(t;X)=1$ and $r\notin B$}
        \State $\widehat G\gets\widehat G\cup\{t\}$; $B\gets B\cup\{r\}$
    \Else
        \State append $(t,\text{violated predicate or cardinality})$ to $J$
    \EndIf
\EndFor
\State \Return $\widehat G,J$
\end{algorithmic}
\end{algorithm}

No gold field value or per-document gold relation-presence indicator is used
by this algorithm. The required document node is externally supplied task
metadata, even though the original benchmark file stores its value inside the
reference record. Our new runner materializes a restricted public-input object
before constructing any prompt, diagnostic, or SHACL shape. Tests change the
reference tails, reference relation set, and defect labels while keeping the
public input fixed and verify that model inputs remain unchanged.

With hash-based schema and duplicate checks, structural normalization takes
expected $O(|G|)$ triple operations. Filtering $k$ candidates costs
$O(k(1+c_X))$, where $c_X$ is the cost of literal source matching, and uses
$O(k+|\mathcal R_X|)$ auxiliary storage. String hashing and comparison still
depend on text length. This bound excludes model inference and does not imply
scalability of a connected-KG reasoning algorithm.

\subsection{A Matched Mechanism Experiment}
\label{sec:impl:enhancement}

The archived Claude experiment compared prompts that differed in instructions
as well as diagnostics. Its system-level advantage cannot identify a unique
mechanism. The new experiment instead uses a single system prompt and input
schema for three generation arms: base input, base plus diagnostic context,
and base plus SHACL validation context. Each response is evaluated both before
and after the identical candidate gate, producing a $3\times2$ comparison.
This pairing isolates the post-generation effect of filtering without another
model call. Comparing diagnostic and base arms estimates the effect of added
diagnostic context under fixed instructions; it does not eliminate generation
variability or establish the necessity of a neural controller.

All new calls use \texttt{google/gemma-4-26B-A4B-it}, temperature zero, and a
4,000-token completion cap. The cap and call budget are matched; input-token
counts differ because diagnostic and SHACL contexts have different lengths.
We report those counts and request wall times, including recorded retries.
Resuming a run does not selectively regenerate failed or malformed responses.

\subsection{Constraint-Aware Literature Baseline}
\label{sec:impl:shacl}

Lin et al.\ propose repair prompts comprising a primer, violation context,
SHACL manifest context, graph context, and instructions
\citep{ref_lin2025}. We implement a document-field adaptation of their full
manifest/full graph ($M+G$) configuration. A real pySHACL validator produces
reports for canonical-head, closed-schema, and maximum-cardinality shapes.
A source-support SPARQL constraint supplies the same literal grounding
information available to our method. Optional fields have no gold-derived
minimum-count constraints.

The adaptation preserves the five-part context design but returns complete
JSON triples instead of SPARQL edits, includes the original source, and groups
violations into one call per document. These changes align the task and budget;
the baseline is labelled \emph{Lin-style SHACL context (adapted)}, not an
exact reproduction of the original Brick/LUBM/QUDT experiments. RDF represents
a set and collapses duplicate triples, while the shared JSON input retains
multiplicity. Both raw and gated outputs are reported so that an extra
validator is not reserved for our arm.

\subsection{Audit of the Earlier Sequential Optimizer}
\label{sec:impl:optimizer_audit}

The repository also contains a more elaborate optimizer with a four-score
quality profile, a learned scope prior, atomic edit bundles, full trial-graph
reassessment, restoration bounds, and utility-based one-step selection.
That optimizer was not called by the archived headline benchmark. We therefore
execute it separately on the frozen candidate responses, holding candidates
fixed and recording every trial graph, profile, utility, acceptance decision,
and stopping reason. Per-relation differences between the normalized input
and proposal define atomic edit bundles; proposal confidence is fixed at 0.7.
The utility/constraint defaults and 12-step horizon are unchanged. We fixed
a source-matching bug that removed whitespace from the source but not from
the candidate value, then repeated the audit with symmetric normalization;
both pre-fix and corrected traces are preserved. We compare
the original learned router with an always-on uniform prior to check whether
routing alone explains any failure.

This execution audit is a test of an alternative selection mechanism. We do
not attribute the one-call pipeline's results to sequential optimization or
reinforcement learning. Its results and limitations are reported alongside
the matched generation and validation experiments.
